\chapter{Guia das Métricas de Avaliação}

Este capítulo apresenta, em linguagem acessível e em definição técnica, cada métrica utilizada na avaliação do subsistema Text-to-SQL do CIC.

\section{Visão geral para o leitor não técnico}

Imagine que você faz uma pergunta ao assistente e ele precisa traduzir sua pergunta em uma ``receita de busca'' no banco de dados (a SQL). Avaliamos essa tradução em camadas:

\begin{itemize}
    \item \textbf{A receita funciona?} (VSR) --- O banco aceita e executa sem erro?
    \item \textbf{A resposta bate com a esperada?} (EX) --- Os números e linhas retornados são os mesmos da referência?
    \item \textbf{Voltou dado quando devia?} (NEA) --- A consulta não ficou vazia sem motivo?
    \item \textbf{Olhou as tabelas certas?} (TSA) --- Usou as tabelas que um especialista usaria?
    \item \textbf{Respeitou os filtros importantes?} (CHS) --- Ano, órgão, tipo de despesa etc. aparecem na SQL?
    \item \textbf{Quanto custou e quanto demorou?} --- Tokens, dólares e segundos de espera.
\end{itemize}

O padrão observado nos experimentos --- VSR alto e EX moderado --- significa: \textit{o modelo quase sempre fala a língua do banco, mas nem sempre interpreta a pergunta exatamente como o especialista que escreveu a referência}.

\section{Execution Accuracy (EX)}

\subsection{Definição leiga}

É a taxa de acerto ``no resultado final'': a SQL gerada devolve \textbf{exatamente} as mesmas linhas e valores que a SQL de referência?

\subsection{Definição técnica}

Dado um par $(q, S_{ref})$ no \textit{golden dataset}, o modelo produz $S_{gen}$. Executam-se ambas no DuckDB, obtendo $R_{ref}$ e $R_{gen}$. Define-se:

\[
\text{EX} = \frac{1}{N} \sum_{i=1}^{N} \mathbb{1}\left[ \text{multiset}(R_{gen}^{(i)}) = \text{multiset}(R_{ref}^{(i)}) \right]
\]

A comparação ignora ordem das linhas, normaliza tipos numéricos (4 casas) e compara colunas por nome (case-insensitive).

\subsection{Por que usamos}

É o padrão-ouro em benchmarks Text-to-SQL (Spider, BIRD). Mede utilidade prática: se o cidadão receberia os mesmos números que a consulta validada.

\subsection{O que inferir}

EX baixo com VSR alto indica \textbf{erro semântico}, não sintático: o modelo consulta o banco corretamente, mas com agregações, filtros ou \textit{joins} diferentes. EX não captura SQLs alternativas semanticamente corretas --- penaliza respostas válidas que divergem da única referência anotada.

\section{Valid SQL Rate (VSR)}

\subsection{Definição leiga}

Quantas vezes a SQL gerada \textbf{roda sem erro} no banco?

\subsection{Definição técnica}

\[
\text{VSR} = \frac{1}{N} \sum_{i=1}^{N} \mathbb{1}\left[ \text{execute}(S_{gen}^{(i)}) \neq \text{erro} \right]
\]

\subsection{Por que usamos}

Distingue falha de sintaxe/tipos (\textit{Binder Error}, coluna inexistente) de falha de intenção. É pré-requisito para qualquer utilidade.

\subsection{O que inferir}

VSR elevado ($>$90\%) confirma que o \textit{prompt} com \\termo{ddl} minificado e amostras é eficaz para ensinar o esquema. Quedas de VSR (ex.: GPT-4o-mini em consultas difíceis) sinalizam limites do modelo menor em \textit{joins} complexos.

\section{Non-Empty Answer (NEA)}

\subsection{Definição leiga}

Quando esperamos dados na resposta, a consulta voltou pelo menos uma linha?

\subsection{Definição técnica}

Para itens com \texttt{espera\_dados=sim}: $\text{NEA} = \mathbb{1}[|R_{gen}| > 0]$. Casos em que a referência retorna vazio são tratados separadamente.

\subsection{Por que usamos}

Evita ``acerto técnico'' com resultado inútil (consulta válida que retorna zero linhas por filtro errado).

\subsection{O que inferir}

Complementa EX: uma SQL pode executar e retornar vazio por ano ou órgão incorretos.

\section{Table Selection Accuracy (TSA)}

\subsection{Definição leiga}

A SQL usou as tabelas que um analista esperaria para aquela pergunta?

\subsection{Definição técnica}

Heurística por expressão regular sobre \texttt{FROM}/\texttt{JOIN}: todas as tabelas listadas em \texttt{tabelas\_esperadas} do golden devem aparecer em $S_{gen}$.

\subsection{Por que usamos}

Erro de tabela é falha grave em bases com dezenas de domínios (frota, folha, licitação). TSA barato computacionalmente.

\subsection{O que inferir}

TSA alto com EX baixo sugere que o modelo ``está no bairro certo, mas na rua errada'' --- tabelas corretas, lógica de filtro/agregação incorreta.

\section{Condition Heuristic Score (CHS)}

\subsection{Definição leiga}

A SQL menciona os filtros importantes (ano, CNPJ, tipo de movimentação)?

\subsection{Definição técnica}

Para condições anotadas $c_1,\ldots,c_k$ em \texttt{condicao\_esperada}:

\[
\text{CHS} = \frac{|\{ j : c_j \subseteq S_{gen} \text{ (substring)} \}|}{k}
\]

\subsection{Por que usamos}

Aproxima ``conformidade com regras de negócio'' sem exigir SQL idêntica.

\subsection{O que inferir}

CHS correlaciona com aderência a convenções municipais (códigos de órgão, categorias econômicas). Útil para diagnóstico por categoria no TCC.

\section{Métricas operacionais}

\subsection{Latência}

\textbf{Geração} ($t_{gen}$): tempo da chamada OpenAI até JSON com \texttt{query\_sql}. \textbf{Execução} ($t_{sql}$): tempo do POST na API DuckDB. Relevante para WhatsApp: usuários toleram poucos segundos.

\subsection{Tokens e custo}

Tokens de entrada incluem o \textit{system prompt} ($\sim$16k--30k). Custo estimado por tabela de preços em \texttt{eval/config.yaml}. Permite comparar Póscagada vs Précagada e justificar modelo em produção.

\section{Síntese interpretativa dos resultados obtidos}

\begin{table}[htbp]
\centering
\caption{Leitura integrada das métricas (experimentos $n=80$)}
\label{tab:leitura-metricas}
\begin{tabular}{|p{3.2cm}|p{9.5cm}|}
\hline
\textbf{Padrão} & \textbf{Interpretação} \\
\hline
VSR alto, EX baixo & Modelo domina sintaxe/esquema; referência e geração divergem semanticamente. \\
\hline
EX sobe com GPT-4o & Modelos maiores melhoram intenção, especialmente em cruzamentos difíceis. \\
\hline
Précagada vs Póscagada & Ganho marginal de EX com custo $\sim$2$\times$; minificação é custo-efetiva. \\
\hline
RH/Folha EX = 0 & Possível efeito do apagão de dados 2020--2022 no contexto + ambiguidade da referência. \\
\hline
\end{tabular}
\end{table}

\section{LLM-as-Judge (próxima etapa)}

As métricas acima não avaliam a qualidade da \textbf{resposta em linguagem natural} ao usuário final. A Seção~\ref{sec:plano-judge} e o documento \texttt{docs/PLANO\_LLM\_JUDGE.md} descrevem o protocolo proposto para julgamento automático por outro \\termo{llm}, ainda não implementado nesta fase.
